\scanpage{180}
Appliquons le th.~3 à la dualité globale, en supposant $X$ \uline{propre et lisse} sur $S$. Soit $E$ module loc. libre sur $X$ (ou complexe parfait $\ldots$), $\widetilde E$ son image inverse sur $X^{p/S}$, \inserted{dont l'image inverse sur $X$ par $f_{X/S}$ est $E^{(p)}$.} Alors on a les homomorphismes
\[
\deleted{\left[H^i(X,E)\xrightarrow{p\text{-lin}}\right]H^i(X^{p/S},\widetilde E)\longrightarrow H^i(X,E^{(p)})}
\]
dit ``puissance p.ième''
\[
 \underset{\substack{\rotatebox{90}{$\simeq$}\\\mathbb R f_*(E)^{(p)}}}{\mathbb R f_*(\widetilde E)}
 \xrightarrow{\unsure{r|\sigma|\varphi}}\mathbb R f_*(E^{(p)})\tag*{($\ast$)}
\]
qui n'est autre essentiellement que l'hom de changemt de base associé au carré commutatif

\diagram{\begin{tikzcd}[ampersand replacement=\&]
{}^{E}X \arrow[r,no head,"F_X"] \arrow[d,no head,"f"] \& X^{F_X^\ast(E)=E^{(p)}}\arrow[d,no head]\\
S\arrow[r,no head,"F_S"]\&S
\end{tikzcd}}

Considérons le \uline{transposé} de $(\ast)$, nous savons que c'est l'hom
\[
\begin{aligned}
&\hspace{10em}\simeq Rg_*\bigl(Rf_{X/S\,*}(E^{(p)}\otimes\Omega^n_{X/S})\bigr)\\
&Rf_*(\check E^{(p)}\otimes\Omega^n_{X/S}[n])\mathrel{\text{---}}
 Rf_*(\check{\widetilde E}\otimes\Omega^n_{\widetilde X/S}[n])
\end{aligned}
\]
associé à l'hom trace
\[
 f_{X/S\,*}(\check E^{(p)}\otimes\Omega^n_{X/S})
 =\check{\widetilde E}\otimes f_{X/S\,*}(\Omega^n_{X/S})
 \longrightarrow\check{\widetilde{\unsure{E|\mathbb E}}}\otimes\Omega^n_{X^{p/S}/S}\tag*{($\ast\ast$)}
\]
déduit, par tensorisation avec $\check{\widetilde E}[n]$ \deleted{\unsure{avec}} de l'hom trace habituel. Celui-ci à son tour s'interprète comme l'hom \deleted{trace} de Cartier.

En particulier, on voit que le transposé de
\[
 \underset{\substack{\rotatebox{90}{$\simeq$}\\R^nf_*(E)^{(p)}}}{R^nf_*(\widetilde E)}
 \longrightarrow R^nf_*(E^{(p)})
\]
\deleted{[supposons $R^nf_*(E)$ et $R^nf_*(E^{(p)})$ loc. libres, pour simplifier]}
% archive sheet number: 179
